% Standalone editable calculation note for issue 15848.
\section{Named Visual Hair Silhouettes}
The four period-inspired styles use the existing CC0 scalp shell, not measured
person-specific geometry.  In the head frame, $x$ points forward, $z$ upward,
and all lengths below are fractions of head length.  A monotone cubic
hairline $z=h(|\theta|)+a\sin(\theta)\max(1-|\theta|/1.4,0)$, with
$\theta=\operatorname{atan2}(y,x-x_c)$ and $|a|\leq 0.15$, cuts
the scalp. The sign of $a$ follows the selected part; zero preserves the
legacy symmetric line. The closed
shell offsets the skin along its normals with the existing crown, fringe,
part and fade thickness fields.  The style presets change the hairline and
these fields; the back hairline ranges from 0.24 to 0.34 head lengths.
The \texttt{side\_swept} and \texttt{swept\_back} styles have fuller crowns
and lower napes; \texttt{slick\_part} and \texttt{receding\_crop} are closer
crops with higher frontal hairlines.  These are visual approximations of
observed grooming, not evidence of likeness or biometric reconstruction.

When any cap is selected, its existing qualified under-hair crop is used
unchanged so crown clearance remains as before. Existing \texttt{short\_crop}
and \texttt{dome} paths remain unchanged. Hair is a non-colliding visual
mesh; the physical spec, body masses, joints and inertias are not changed.
Acceptance requires closed outward finite meshes, distinct silhouettes,
legacy and cap geometry equality, and matching native MuJoCo dynamics at
representative states. Reproduce with
\begin{verbatim}
PYTHONPATH=.:vendor/ud-tools/src MUJOCO_GL=egl \
  python -m pytest tests/unit/model_appearance/test_hair_styles.py
PYTHONPATH=.:vendor/ud-tools/src MUJOCO_GL=egl \
  python scripts/render_hair_styles.py --out /tmp/hair-styles
\end{verbatim}
